\chapter{References}
\label{ch:references}

\begin{enumerate}[label={[\arabic*]}, leftmargin=2.4em, itemsep=0.55em]

\item\label{burns2016borg} B. Burns, B. Grant, D. Oppenheimer, E. Brewer, and J. Wilkes, ``Borg, Omega, and Kubernetes,'' \textit{Commun. ACM}, vol. 59, no. 5, pp. 50--57, May 2016.

\item\label{lorido2014review} T. Lorido-Botran, J. Miguel-Alonso, and J. A. Lozano, ``A review of auto-scaling techniques for elastic applications in cloud environments,'' \textit{J. Grid Comput.}, vol. 12, no. 4, pp. 559--592, Dec. 2014.

\item\label{qu2018survey} C. Qu, R. N. Calheiros, and R. Buyya, ``Auto-scaling web applications in clouds: A taxonomy and survey,'' \textit{ACM Comput. Surv.}, vol. 51, no. 4, pp. 73:1--73:33, Jul. 2018.

\item\label{k8s_hpa_docs} The Kubernetes Authors, ``Horizontal Pod Autoscaling,'' Kubernetes Documentation. [Online]. Available: \url{https://kubernetes.io/docs/tasks/run-application/horizontal-pod-autoscale/}

\item\label{sutton2018rl} R. S. Sutton and A. G. Barto, \textit{Reinforcement Learning: An Introduction}, 2nd ed. Cambridge, MA: MIT Press, 2018.

\item\label{mao2016resource} H. Mao, M. Alizadeh, I. Menache, and S. Kandula, ``Resource management with deep reinforcement learning,'' in \textit{Proc. 15th ACM Workshop Hot Topics Netw. (HotNets)}, 2016, pp. 50--56.

\item\label{horovitz2018qlearning} S. Horovitz and Y. Arian, ``Efficient cloud auto-scaling with SLA objective using Q-learning,'' in \textit{Proc. IEEE 6th Int. Conf. Future Internet Things Cloud (FiCloud)}, 2018, pp. 85--92.

\item\label{rossi2019horizontal} F. Rossi, M. Nardelli, and V. Cardellini, ``Horizontal and vertical scaling of container-based applications using reinforcement learning,'' in \textit{Proc. IEEE 12th Int. Conf. Cloud Comput. (CLOUD)}, 2019, pp. 329--338.

\item\label{imdoukh2020lstm} M. Imdoukh, I. Ahmad, and M. G. Alfailakawi, ``Machine learning-based auto-scaling for containerized applications,'' \textit{Neural Comput. Appl.}, vol. 32, pp. 9745--9760, 2020.

\item\label{schulman2017ppo} J. Schulman, F. Wolski, P. Dhariwal, A. Radford, and O. Klimov, ``Proximal policy optimization algorithms,'' 2017, arXiv:1707.06347. [Online]. Available: \url{https://arxiv.org/abs/1707.06347}

\item\label{towers2024gymnasium} M. Towers et al., ``Gymnasium: A standard interface for reinforcement learning environments,'' 2024, arXiv:2407.17032. [Online]. Available: \url{https://arxiv.org/abs/2407.17032}

\item\label{taherizadeh2019dynamic} S. Taherizadeh and V. Stankovski, ``Dynamic multi-level auto-scaling rules for containerized applications,'' \textit{Comput. J.}, vol. 62, no. 2, pp. 174--197, Feb. 2019.

\item\label{rzadca2020autopilot} K. Rzadca et al., ``Autopilot: Workload autoscaling at Google,'' in \textit{Proc. 15th Eur. Conf. Comput. Syst. (EuroSys)}, 2020, pp. 16:1--16:16.

\item\label{rattihalli2019vertical} G. Rattihalli, M. Govindaraju, H. Lu, and D. Tiwari, ``Exploring potential for non-disruptive vertical auto scaling and resource estimation in Kubernetes,'' in \textit{Proc. IEEE 12th Int. Conf. Cloud Comput. (CLOUD)}, 2019, pp. 33--40.

\item\label{toka2021edge} L. Toka, G. Dobreff, B. Fodor, and B. Sonkoly, ``Machine learning-based scaling management for Kubernetes edge clusters,'' \textit{IEEE Trans. Netw. Serv. Manag.}, vol. 18, no. 1, pp. 958--972, Mar. 2021.

\item\label{qiu2020firm} H. Qiu, S. S. Banerjee, S. Jha, Z. T. Kalbarczyk, and R. K. Iyer, ``FIRM: An intelligent fine-grained resource management framework for SLO-oriented microservices,'' in \textit{Proc. 14th USENIX Symp. Oper. Syst. Des. Implement. (OSDI)}, 2020, pp. 805--825.

\item\label{henderson2018matters} P. Henderson, R. Islam, P. Bachman, J. Pineau, D. Precup, and D. Meger, ``Deep reinforcement learning that matters,'' in \textit{Proc. 32nd AAAI Conf. Artif. Intell.}, 2018, pp. 3207--3214.

\item\label{agarwal2021precipice} R. Agarwal, M. Schwarzer, P. S. Castro, A. C. Courville, and M. G. Bellemare, ``Deep reinforcement learning at the edge of the statistical precipice,'' in \textit{Proc. Adv. Neural Inf. Process. Syst. (NeurIPS)}, vol. 34, 2021, pp. 29304--29320.

\item\label{dulacarnold2019challenges} G. Dulac-Arnold, D. Mankowitz, and T. Hester, ``Challenges of real-world reinforcement learning,'' 2019, arXiv:1904.12901. [Online]. Available: \url{https://arxiv.org/abs/1904.12901}

\item\label{garcia2015safe} J. Garc\'ia and F. Fern\'andez, ``A comprehensive survey on safe reinforcement learning,'' \textit{J. Mach. Learn. Res.}, vol. 16, no. 42, pp. 1437--1480, 2015.

\item\label{nosek2018preregistration} B. A. Nosek, C. R. Ebersole, A. C. DeHaven, and D. T. Mellor, ``The preregistration revolution,'' \textit{Proc. Natl. Acad. Sci. USA}, vol. 115, no. 11, pp. 2600--2606, Mar. 2018.

\item\label{kaelbling1998pomdp} L. P. Kaelbling, M. L. Littman, and A. R. Cassandra, ``Planning and acting in partially observable stochastic domains,'' \textit{Artif. Intell.}, vol. 101, no. 1--2, pp. 99--134, May 1998.

\item\label{schulman2016gae} J. Schulman, P. Moritz, S. Levine, M. Jordan, and P. Abbeel, ``High-dimensional continuous control using generalized advantage estimation,'' in \textit{Proc. Int. Conf. Learn. Represent. (ICLR)}, 2016. [Online]. Available: \url{https://arxiv.org/abs/1506.02438}

\item\label{raffin2021sb3} A. Raffin, A. Hill, A. Gleave, A. Kanervisto, M. Ernestus, and N. Dormann, ``Stable-Baselines3: Reliable reinforcement learning implementations,'' \textit{J. Mach. Learn. Res.}, vol. 22, no. 268, pp. 1--8, 2021.

\item\label{prometheus_docs} The Prometheus Authors, ``Prometheus Documentation.'' [Online]. Available: \url{https://prometheus.io/docs/introduction/overview/}

\item\label{minikube_docs} The Kubernetes Authors, ``Minikube Documentation.'' [Online]. Available: \url{https://minikube.sigs.k8s.io/docs/}

\end{enumerate}
